% ============================================================================
% CHAPTER 5: OPERATIONS & MAINTENANCE
% ============================================================================

\chapter{Operations \& Maintenance}

\section{Starting and Stopping Services}

\subsection{Complete System Startup}

\begin{commandbox}
\begin{lstlisting}[style=bashstyle]
# Navigate to project root
cd /Users/sreehasgopinathan/Documents/Auburn/Project/Data_Ingestion

# Start all services
docker-compose up -d

# Verify services are healthy (wait 30-60 seconds)
docker-compose ps

# Check logs for any errors
docker-compose logs -f --tail=100
\end{lstlisting}
\end{commandbox}

\textbf{Expected Startup Time:}
\begin{itemize}
    \item Database services: 15-20 seconds
    \item Backend services: 20-30 seconds
    \item Frontend: 10-15 seconds
    \item Total: ~60 seconds for full stack
\end{itemize}

\subsection{Starting Individual Services}

\begin{commandbox}
\begin{lstlisting}[style=bashstyle]
# Start only database services
docker-compose up -d postgres elasticsearch neo4j redis

# Start only application services
docker-compose up -d api-gateway pv-classification-service frontend

# Start specific service
docker-compose up -d frontend
\end{lstlisting}
\end{commandbox}

\subsection{Graceful Shutdown}

\begin{commandbox}
\begin{lstlisting}[style=bashstyle]
# Stop all services (graceful shutdown, 10s timeout)
docker-compose stop

# Stop all services immediately (force)
docker-compose down

# Stop and remove volumes (WARNING: Deletes all data)
docker-compose down -v
\end{lstlisting}
\end{commandbox}

\begin{warningbox}
\textbf{Data Loss Warning:} Using \texttt{docker-compose down -v} will delete all database data, uploaded documents, and configuration. Only use for complete system reset.
\end{warningbox}

\subsection{Service Restart}

\begin{commandbox}
\begin{lstlisting}[style=bashstyle]
# Restart specific service (e.g., after config change)
docker-compose restart frontend

# Restart all services
docker-compose restart

# Rebuild and restart (after code changes)
docker-compose up -d --build frontend
\end{lstlisting}
\end{commandbox}

\section{Health Monitoring}

\subsection{Service Health Checks}

\textbf{Docker Health Status:}
\begin{commandbox}
\begin{lstlisting}[style=bashstyle]
# View service status
docker-compose ps

# Expected output:
# NAME                    STATUS              PORTS
# frontend                Up (healthy)        0.0.0.0:3000->80/tcp
# api-gateway             Up (healthy)        0.0.0.0:8000->8000/tcp
# pv-classification       Up (healthy)        0.0.0.0:8012->8002/tcp
# postgres                Up (healthy)        0.0.0.0:5432->5432/tcp
# elasticsearch           Up (healthy)        0.0.0.0:9200->9200/tcp
# neo4j                   Up (healthy)        0.0.0.0:7474->7474/tcp
\end{lstlisting}
\end{commandbox}

\textbf{Manual Health Check Endpoints:}
\begin{commandbox}
\begin{lstlisting}[style=bashstyle]
# API Gateway
curl http://localhost:8000/health

# Expected: {"status":"healthy","version":"1.0.0"}

# Elasticsearch
curl http://localhost:9200/_cluster/health

# Neo4j (via Cypher)
curl -u neo4j:password http://localhost:7474/db/data/
\end{lstlisting}
\end{commandbox}

\subsection{Log Monitoring}

\textbf{View Real-time Logs:}
\begin{commandbox}
\begin{lstlisting}[style=bashstyle]
# All services
docker-compose logs -f

# Specific service
docker-compose logs -f api-gateway

# Last 100 lines
docker-compose logs --tail=100 api-gateway

# With timestamps
docker-compose logs -f --timestamps api-gateway

# Filter by pattern (e.g., errors)
docker-compose logs -f api-gateway | grep ERROR
\end{lstlisting}
\end{commandbox}

\textbf{Log Locations (Inside Containers):}
\begin{table}[h]
\centering
\begin{tabularx}{\textwidth}{lX}
\toprule
\textbf{Service} & \textbf{Log Path} \\
\midrule
API Gateway & /app/logs/api-gateway.log \\
PV Classification & /app/logs/pv-classification.log \\
Frontend (Nginx) & /var/log/nginx/access.log, error.log \\
PostgreSQL & /var/log/postgresql/ \\
\bottomrule
\end{tabularx}
\end{table}

\subsection{Prometheus \& Grafana Monitoring}

\textbf{Access Dashboards:}
\begin{itemize}
    \item Prometheus: \texttt{http://localhost:9090}
    \item Grafana: \texttt{http://localhost:3001} (admin/admin)
\end{itemize}

\textbf{Key Metrics to Monitor:}
\begin{enumerate}
    \item \textbf{System Resources}
    \begin{itemize}
        \item CPU usage per container
        \item Memory usage and limits
        \item Disk I/O and space
    \end{itemize}
    
    \item \textbf{Application Metrics}
    \begin{itemize}
        \item Request rate (requests/second)
        \item Response time (p50, p95, p99)
        \item Error rate (4xx, 5xx responses)
    \end{itemize}
    
    \item \textbf{Database Metrics}
    \begin{itemize}
        \item Query execution time
        \item Connection pool usage
        \item Index hit rate
    \end{itemize}
    
    \item \textbf{Queue Metrics}
    \begin{itemize}
        \item Processing queue depth
        \item Job success/failure rate
        \item Average processing time
    \end{itemize}
\end{enumerate}

\begin{figure}[h]
\centering
\fbox{\textit{[SCREENSHOT PLACEHOLDER: Grafana Dashboard - System Overview]}}
\caption{Grafana Monitoring Dashboard}
\label{fig:grafana-monitoring}
\end{figure}

\section{Database Maintenance}

\subsection{PostgreSQL Maintenance}

\textbf{Backup Database:}
\begin{commandbox}
\begin{lstlisting}[style=bashstyle]
# Full database backup
docker-compose exec postgres pg_dump -U postgres document_intelligence > backup_$(date +%Y%m%d).sql

# Backup with compression
docker-compose exec postgres pg_dump -U postgres document_intelligence | gzip > backup_$(date +%Y%m%d).sql.gz

# Restore from backup
docker-compose exec -T postgres psql -U postgres document_intelligence < backup_20241101.sql
\end{lstlisting}
\end{commandbox}

\textbf{Database Vacuum (Performance Optimization):}
\begin{commandbox}
\begin{lstlisting}[style=bashstyle]
# Analyze and optimize tables
docker-compose exec postgres psql -U postgres -d document_intelligence -c "VACUUM ANALYZE;"

# Full vacuum (requires downtime)
docker-compose exec postgres psql -U postgres -d document_intelligence -c "VACUUM FULL;"

# Check database size
docker-compose exec postgres psql -U postgres -c "\l+ document_intelligence"
\end{lstlisting}
\end{commandbox}

\subsection{Elasticsearch Maintenance}

\textbf{Index Management:}
\begin{commandbox}
\begin{lstlisting}[style=bashstyle]
# Check cluster health
curl http://localhost:9200/_cluster/health?pretty

# List all indices
curl http://localhost:9200/_cat/indices?v

# Check index size
curl http://localhost:9200/_cat/indices/documents?v&h=index,docs.count,store.size

# Force merge (optimize) index
curl -X POST "http://localhost:9200/documents/_forcemerge?max_num_segments=1"
\end{lstlisting}
\end{commandbox}

\textbf{Snapshot Backup:}
\begin{commandbox}
\begin{lstlisting}[style=bashstyle]
# Register snapshot repository
curl -X PUT "http://localhost:9200/_snapshot/backup_repo" -H 'Content-Type: application/json' -d'
{
  "type": "fs",
  "settings": {
    "location": "/usr/share/elasticsearch/backups"
  }
}'

# Create snapshot
curl -X PUT "http://localhost:9200/_snapshot/backup_repo/snapshot_$(date +%Y%m%d)?wait_for_completion=true"

# Restore snapshot
curl -X POST "http://localhost:9200/_snapshot/backup_repo/snapshot_20241101/_restore"
\end{lstlisting}
\end{commandbox}

\subsection{Neo4j Maintenance}

\textbf{Database Backup:}
\begin{commandbox}
\begin{lstlisting}[style=bashstyle]
# Stop Neo4j for consistent backup
docker-compose stop neo4j

# Backup database files
docker run --rm -v data_ingestion_neo4j-data:/data -v $(pwd):/backup \
  ubuntu tar czf /backup/neo4j_backup_$(date +%Y%m%d).tar.gz /data

# Restart Neo4j
docker-compose start neo4j
\end{lstlisting}
\end{commandbox}

\textbf{Database Statistics:}
\begin{commandbox}
\begin{lstlisting}[style=bashstyle]
# Connect to Neo4j browser (http://localhost:7474)
# Run Cypher query:

// Count nodes by type
MATCH (n)
RETURN labels(n) as NodeType, count(*) as Count
ORDER BY Count DESC;

// Count relationships by type
MATCH ()-[r]->()
RETURN type(r) as RelationshipType, count(*) as Count
ORDER BY Count DESC;

// Database size
CALL apoc.meta.stats();
\end{lstlisting}
\end{commandbox}

\section{Performance Optimization}

\subsection{Docker Resource Allocation}

\textbf{Recommended Resource Limits:}

\begin{lstlisting}[style=pythonstyle]
# docker-compose.yml
services:
  elasticsearch:
    mem_limit: 4g
    cpus: 2.0
    environment:
      - "ES_JAVA_OPTS=-Xms2g -Xmx2g"
  
  postgres:
    mem_limit: 2g
    cpus: 1.0
  
  neo4j:
    mem_limit: 2g
    cpus: 1.0
    environment:
      - NEO4J_dbms_memory_heap_max__size=1g
  
  pv-classification-service:
    mem_limit: 4g  # For BERT model
    cpus: 2.0
\end{lstlisting}

\subsection{Database Query Optimization}

\textbf{PostgreSQL Indexing:}
\begin{commandbox}
\begin{lstlisting}[style=bashstyle]
# Create indexes for frequently queried fields
docker-compose exec postgres psql -U postgres -d document_intelligence -c "
CREATE INDEX IF NOT EXISTS idx_documents_space_id ON documents(space_id);
CREATE INDEX IF NOT EXISTS idx_documents_classification ON documents(classification);
CREATE INDEX IF NOT EXISTS idx_documents_created_at ON documents(created_at DESC);
CREATE INDEX IF NOT EXISTS idx_documents_doc_id ON documents(doc_id);
"
\end{lstlisting}
\end{commandbox}

\textbf{Elasticsearch Query Optimization:}
\begin{itemize}
    \item Use \texttt{\_source} filtering to reduce payload size
    \item Implement query result caching (Redis)
    \item Use \texttt{scroll} API for large result sets
    \item Disable scoring when not needed (\texttt{filter} context)
\end{itemize}

\subsection{Caching Strategy}

\textbf{Redis Caching:}
\begin{lstlisting}[style=pythonstyle]
# Classification result caching (30 minutes TTL)
CACHE_KEY = f"classification:{doc_id}"
cached_result = redis.get(CACHE_KEY)

if cached_result:
    return json.loads(cached_result)

# Compute result
result = classify_document(document)

# Cache result
redis.setex(CACHE_KEY, 1800, json.dumps(result))
return result
\end{lstlisting}

\textbf{Frontend Caching:}
\begin{itemize}
    \item API responses cached in React Query (5 minutes)
    \item Static assets cached with aggressive headers (1 year)
    \item Service Worker for offline capability (future)
\end{itemize}

\section{Security Operations}

\subsection{SSL/TLS Certificate Management}

\textbf{Generate Self-Signed Certificate (Development):}
\begin{commandbox}
\begin{lstlisting}[style=bashstyle]
# Generate certificate
openssl req -x509 -nodes -days 365 -newkey rsa:2048 \
  -keyout ./nginx/ssl/selfsigned.key \
  -out ./nginx/ssl/selfsigned.crt \
  -subj "/C=US/ST=Alabama/L=Auburn/O=Auburn/CN=localhost"

# Update nginx.conf to use SSL
# Restart frontend
docker-compose restart frontend
\end{lstlisting}
\end{commandbox}

\textbf{Let's Encrypt Certificate (Production):}
\begin{commandbox}
\begin{lstlisting}[style=bashstyle]
# Install certbot
sudo apt-get install certbot python3-certbot-nginx

# Obtain certificate
sudo certbot --nginx -d yourdomain.com -d www.yourdomain.com

# Auto-renewal (cron job)
0 0 * * * certbot renew --quiet --post-hook "docker-compose restart frontend"
\end{lstlisting}
\end{commandbox}

\subsection{Secret Rotation}

\textbf{JWT Secret Rotation:}
\begin{commandbox}
\begin{lstlisting}[style=bashstyle]
# Generate new secret
NEW_SECRET=$(openssl rand -hex 32)

# Update .env file
sed -i "s/JWT_SECRET_KEY=.*/JWT_SECRET_KEY=$NEW_SECRET/" .env

# Restart API Gateway (invalidates all existing tokens)
docker-compose restart api-gateway

# Notify users to re-login
\end{lstlisting}
\end{commandbox}

\textbf{Database Password Rotation:}
\begin{commandbox}
\begin{lstlisting}[style=bashstyle]
# Generate new password
NEW_PASSWORD=$(openssl rand -base64 24)

# Update PostgreSQL password
docker-compose exec postgres psql -U postgres -c \
  "ALTER USER postgres WITH PASSWORD '$NEW_PASSWORD';"

# Update .env file
sed -i "s/POSTGRES_PASSWORD=.*/POSTGRES_PASSWORD=$NEW_PASSWORD/" .env

# Restart services that connect to database
docker-compose restart api-gateway pv-classification-service
\end{lstlisting}
\end{commandbox}

\subsection{Audit Log Review}

\textbf{Security Event Monitoring:}
\begin{commandbox}
\begin{lstlisting}[style=bashstyle]
# View failed login attempts
docker-compose logs api-gateway | grep "Failed login"

# View unauthorized access attempts
docker-compose logs api-gateway | grep "401\|403"

# Export audit logs for compliance
docker-compose exec postgres psql -U postgres -d document_intelligence -c \
  "COPY (SELECT * FROM audit_logs WHERE created_at >= NOW() - INTERVAL '30 days') \
   TO STDOUT CSV HEADER" > audit_export_$(date +%Y%m%d).csv
\end{lstlisting}
\end{commandbox}

\section{Scaling Operations}

\subsection{Horizontal Scaling}

\textbf{Scale API Gateway:}
\begin{commandbox}
\begin{lstlisting}[style=bashstyle]
# Scale to 3 instances
docker-compose up -d --scale api-gateway=3

# Add load balancer (nginx)
# Edit docker-compose.yml to add nginx load balancer service
\end{lstlisting}
\end{commandbox}

\textbf{Load Balancer Configuration (nginx):}
\begin{lstlisting}[style=pythonstyle]
upstream api_backend {
    least_conn;  # Load balancing algorithm
    server api-gateway-1:8000 max_fails=3 fail_timeout=30s;
    server api-gateway-2:8000 max_fails=3 fail_timeout=30s;
    server api-gateway-3:8000 max_fails=3 fail_timeout=30s;
}

server {
    listen 80;
    location /api/ {
        proxy_pass http://api_backend;
        proxy_set_header Host $host;
        proxy_set_header X-Real-IP $remote_addr;
    }
}
\end{lstlisting}

\subsection{Vertical Scaling}

\textbf{Increase Resource Limits:}
\begin{lstlisting}[style=pythonstyle]
# docker-compose.yml
services:
  pv-classification-service:
    deploy:
      resources:
        limits:
          cpus: '4.0'      # Increase from 2.0
          memory: 8G       # Increase from 4G
        reservations:
          cpus: '2.0'
          memory: 4G
\end{lstlisting}

\begin{commandbox}
\begin{lstlisting}[style=bashstyle]
# Apply changes
docker-compose up -d --force-recreate pv-classification-service
\end{lstlisting}
\end{commandbox}

\subsection{Database Sharding (Advanced)}

\textbf{Elasticsearch Sharding Configuration:}
\begin{commandbox}
\begin{lstlisting}[style=bashstyle]
# Update index settings
curl -X PUT "http://localhost:9200/documents/_settings" -H 'Content-Type: application/json' -d'
{
  "index": {
    "number_of_shards": 5,
    "number_of_replicas": 1
  }
}'
\end{lstlisting}
\end{commandbox}

\section{Disaster Recovery}

\subsection{Backup Strategy}

\textbf{Recommended Backup Schedule:}
\begin{table}[h]
\centering
\caption{Backup Frequency Recommendations}
\begin{tabularx}{\textwidth}{lXl}
\toprule
\textbf{Data Type} & \textbf{Backup Frequency} & \textbf{Retention} \\
\midrule
PostgreSQL & Daily (full), Hourly (incremental) & 30 days \\
Elasticsearch & Daily snapshots & 14 days \\
Neo4j & Weekly & 30 days \\
Uploaded Files & Daily (rsync) & 90 days \\
Configuration & After each change & Indefinite (Git) \\
\bottomrule
\end{tabularx}
\end{table}

\textbf{Automated Backup Script:}
\begin{commandbox}
\begin{lstlisting}[style=bashstyle]
#!/bin/bash
# daily_backup.sh

BACKUP_DIR="/backups/$(date +%Y%m%d)"
mkdir -p $BACKUP_DIR

# PostgreSQL backup
docker-compose exec -T postgres pg_dump -U postgres document_intelligence \
  | gzip > $BACKUP_DIR/postgres.sql.gz

# Elasticsearch snapshot
curl -X PUT "http://localhost:9200/_snapshot/backup_repo/snapshot_$(date +%Y%m%d)?wait_for_completion=true"

# Copy Neo4j data
docker run --rm -v data_ingestion_neo4j-data:/data -v $BACKUP_DIR:/backup \
  ubuntu tar czf /backup/neo4j.tar.gz /data

# Upload to S3 (optional)
aws s3 sync $BACKUP_DIR s3://your-backup-bucket/$(date +%Y%m%d)/

# Clean old backups (keep last 30 days)
find /backups/* -type d -mtime +30 -exec rm -rf {} \;
\end{lstlisting}
\end{commandbox}

\textbf{Schedule with Cron:}
\begin{commandbox}
\begin{lstlisting}[style=bashstyle]
# Edit crontab
crontab -e

# Add daily backup at 2 AM
0 2 * * * /path/to/daily_backup.sh >> /var/log/backup.log 2>&1
\end{lstlisting}
\end{commandbox}

\subsection{Disaster Recovery Procedure}

\textbf{Full System Recovery:}
\begin{enumerate}
    \item Stop all services: \texttt{docker-compose down}
    \item Restore database backups (PostgreSQL, Elasticsearch, Neo4j)
    \item Restore uploaded files and configuration
    \item Start services: \texttt{docker-compose up -d}
    \item Verify data integrity
    \item Update DNS/load balancer to point to recovered system
\end{enumerate}

\textbf{Recovery Time Objective (RTO):} 4 hours

\textbf{Recovery Point Objective (RPO):} 24 hours (last daily backup)

\end{document}

